%! Characteristics of Polynomial Functions I — Unit 3, Lesson 3.4 — Extra Credit (Answer Key)
% MIRROR of ../extra_credit/main.tex: identical except -boxes -> -key and the page-header title.
% UNITS 3+ PACKET-FIRST shape (LESSON_SHAPE.md §0a, 2026-11-14). GENERATED, GRADED, OPTIONAL.
% Last in the packet; the build starts it on a fresh sheet so it tears off. TWO PAGES, exactly,
% blank and key. Five stretch items, harder than every ★★★:
%   q1 COMBINE  — degree, leading coefficient and end behavior of a three-factor product
%   q2 REVERSE  — find k from a stated leading coefficient; find c that makes a function even
%   q3 MODEL    — a campus-garden harvest model, its end behavior, and why the model must stop
%   q4 FIND THE MISTAKE — "x^3 + 1 is odd because its degree is odd" (rubric)
%   q5 SOL-style MULTIPLE CHOICE — the crux distractor (C) reads the first-written term
% Contexts (garden) differ from the notes' (egg drop, spring formal): teach, then transfer.
% Every \gradedblank line is copied VERBATIM into the key. NO NAME LINE.
\documentclass[12pt]{article}
\usepackage{algebra2-article}
\usepackage{algebra2-key}
\usepackage{grading}          % math-shared: key mode prints each answer
\gradingassignment{alg2-u03-l04-ec}
\usepackage{algebra2-graded}
\gradingcourse{algebra_2}

\begin{document}

\pageheader{Unit 3, Lesson 3.4}{Extra Credit --- Answer Key}

\vspace{0.10in}

\boxguard[8]
\begin{remindbox}
\small
\textbf{Extra credit --- optional.} Tear off these pages and turn them in with your homework,
the first class after two study halls. I will announce the due date in class and on TurtleNet.
{\small\itshape Write each final answer in its box. Only what is inside a box is graded.}\par
\end{remindbox}

\vspace{0.08in}

\boxguard[20]
\begin{notesbox}{Stretch}
\small
\begin{enumerate}[leftmargin=1.9em, itemsep=10pt]
  \item \textbf{Combine.} Let $f(x)=-x(2x-1)^2(x+3)$. Find the degree and the leading
    coefficient, then describe each end of the graph as \emph{up} or \emph{down}.
    \begin{work}
      \text{degree} &= 1+2+1=4 && \text{add the degrees} \\
      \text{leading coefficient} &= (-1)(2)^2(1)=-4 && \text{$(2x-1)^2$ leads with $4x^2$} \\
      &\Rightarrow \text{even and negative: both ends fall}
    \end{work}
    (a) Degree: \inblank{\blank{2.0cm}}{\gradedblank[id=q1a, type=short, points=1, answer={4}, check=numeric]{2.0cm-2pt}{12.5pt}}
    \qquad (b) Leading coefficient: \inblank{\blank{2.0cm}}{\gradedblank[id=q1b, type=short, points=1, answer={-4}, check=numeric]{2.0cm-2pt}{12.5pt}}

    \vspace{6pt}
    (c) Left end: \inblank{\blank{2.0cm}}{\gradedblank[id=q1c, type=short, points=1, answer={down}, check=text]{2.0cm-2pt}{12.5pt}}
    \qquad (d) Right end: \inblank{\blank{2.0cm}}{\gradedblank[id=q1d, type=short, points=1, answer={down}, check=text]{2.0cm-2pt}{12.5pt}}

  \item \textbf{Reverse it.}
    (a) The polynomial $p(x)=(kx+1)(x-2)^2(3-x)$ has leading coefficient $12$. Find $k$.
    \begin{work}
      \text{leading coefficient} &= k\cdot1\cdot(-1)=-k && \text{the $x$-terms are $kx$, $x^2$, $-x$} \\
      -k &= 12 \\
      k &= -12
    \end{work}
    $k=$ \inblank{\blank{2.0cm}}{\gradedblank[id=q2a, type=short, points=2, answer={-12}, check=numeric]{2.0cm-2pt}{12.5pt}}

    \vspace{6pt}
    (b) Find the value of $c$ that makes $h(x)=x^4+(c-3)x^3+5$ an \emph{even} function.
    \begin{work}
      \text{even needs only even powers} &\Rightarrow \text{the $x^3$ term must vanish} \\
      c-3 &= 0 \\
      c &= 3
    \end{work}
    $c=$ \inblank{\blank{2.0cm}}{\gradedblank[id=q2b, type=short, points=2, answer={3}, check=numeric]{2.0cm-2pt}{12.5pt}}

  \item \textbf{Model it.} The garden club models the number of tomatoes harvested $x$ weeks
    after planting by $T(x)=30x^2-2x^3$, for $0\le x\le15$.
    \begin{work}
      \text{standard form: } T(x) &= -2x^3+30x^2 && \text{leading term } -2x^3
    \end{work}
    (a) Degree: \inblank{\blank{2.0cm}}{\gradedblank[id=q3a, type=short, points=1, answer={3}, check=numeric]{2.0cm-2pt}{12.5pt}}
    \qquad (b) Leading coefficient: \inblank{\blank{2.0cm}}{\gradedblank[id=q3b, type=short, points=1, answer={-2}, check=numeric]{2.0cm-2pt}{12.5pt}}

    \vspace{6pt}
    (c) Write the right-end behavior of $T$ in arrow notation. What would the model predict for
    the harvest many weeks later, and why must the club stop using it after week $15$?
    \rubricbox{\gradedblank[id=q3c, type=rubric, points=2, display={As $x\to\infty$, $T(x)\to-\infty$: the model would predict a negative number of tomatoes, which is impossible. $T(15)=6750-6750=0$, so after week $15$ it gives negative counts and no longer fits.}]{\linewidth}{2in}}
\end{enumerate}
\end{notesbox}

\boxguard[24]
\begin{notesbox}{Stretch, continued}
\small
\begin{enumerate}[leftmargin=1.9em, itemsep=10pt, start=4]
  \item \textbf{Find the mistake.} Student A classified $g(x)=x^3+1$:
    \begin{center}
    \begin{tabular}{@{}r@{\;}l@{}}
      1. & The degree of $g$ is $3$, which is odd. \\
      2. & So $g$ is an odd function, and its graph is symmetric about the origin. \\
    \end{tabular}
    \end{center}
    Which step is wrong? Explain the error, show the correct test, and classify $g$.
    \rubricbox{\gradedblank[id=q4, type=rubric, points=3, display={Step 2. An odd degree does not make an odd function; test $g(-x)$. $g(-x)=-x^3+1$, which is neither $g(x)=x^3+1$ nor $-g(x)=-x^3-1$, so $g$ is neither (the constant $1$ is an even power, $x^0$).}]{\linewidth}{2in}}

  \item \textbf{Multiple choice (SOL-style).} Which function has this end behavior: as
    $x\to-\infty$, $f(x)\to-\infty$, and as $x\to\infty$, $f(x)\to-\infty$? Write the letter of
    the single best answer in its Letter box.
    \begin{enumerate}[label=(\Alph*), leftmargin=2.0em, itemsep=1pt]
      \item $f(x)=4x^3-x^2$
      \item $f(x)=x^2-3x^4$
      \item $f(x)=-2x+x^4$
      \item $f(x)=-x^5+x^2$
    \end{enumerate}
    \letterbox{\gradedblank[id=q5, type=choice, points=1, answer={B}]{0.4in}{0.4in}}
\end{enumerate}
\end{notesbox}

\end{document}
